\section{An Introduction to the Campaign}

\begin{multicols}{2}

Welcome. Our Dragon Quest (DQ) campaign started in 1981 with a single
game master (GM), Robert Leyland.

Since that time and with the help of more than a few GMs the campaign
has expanded to encompass roughly 50–60 players, over 100 active
characters, roughly 10 hard working GMs and approximately as many
worlds as 15 GMs can dream up.

\subsection{How It Works}

Guild meetings are held four times a year, in the real world and the
game world, for players and GMs to arrange games and swap stories.  At
the meeting GMs announce the games that they will be running the
following season and the players sign up. There are in character
adventure recaps, awards for bravest, smartest etc (story time) and
announcements for characters.

Our in character newspaper, “The Seagate Times”, is also available at
the meeting.  Finally there are out of character announcements of rule
changes etc, and then we adjourn to a cafe.

The GMs meet half an hour prior to each Guild meeting to discuss
campaign issues, resolve rules problems and generally give each other
support.

The meetings are held on the second Sundays of March, June, September
and December. They start at 12.30 pm for GMs and 1pm for players and
finish at about 3pm.  They are held at a variety of locations so
please check the DQ Wiki.

\subsection{Games}

Games run for three months, usually for 3½ hours one night per
week. The adventure will take less than that three months in the game
world, and the remainder of the time can be spent in training, ready
to go out the next season, when real world and game time are
synchronised on a one to one basis.  Players are expected to be
available for all sessions of a game. If you want to play but think
you might have to miss more than a session or two, discuss this in
advance with your GM.

\subsection{Advantages}

\begin{itemize}

\item Longevity – this campaign has been going on for 41 years (as at
  2022).
\item If you don’t like the group you are playing with, you can play
  with different people next season.
\item If your GM runs out of steam, you can still play your character.
\item You can take a break and come back.
\item You don’t ever have to take a break
\item You only need to commit for 13 weeks at a time.

\end{itemize}

\subsection{After the Adventure}

\subsubsection{Experience Points}

At the end of each season (after three months of play) characters are
rewarded with Experience Points (EP), money and loot.  The character
spends EP training to improve abilities or learn new ones. This is
called ‘ranking’. EP is awarded for attendance, role-playing,
contribution and difficulty level.

\begin{description}
\item{Attendance} turning up on time, with character sheets completed
  and dice, paper etc. to hand.
\item{Role-playing} skill at role-playing and consistently staying in
  character through out the game.
\item{Contribution} partly how much your character does, but primarily
  how you as a player contribute to the scenario and to everybody’s
  enjoyment of the game.
\item{Difficulty} the difficulty level is announced at the start of
  the game and is an indication of the danger and the level of ability
  required to successfully complete the scenario.
\end{description}

\subsubsection{Treasure Splits}

\begin{itemize}
\item In most circumstances any treasure gained on adventure is valued
  and split equally amongst the party members. Where there are items
  to be shared out then party members roll dice and, in the order
  lowest roll to highest, each person chooses one item until all have
  chosen. If there are still items to be shared out then everyone gets
  a second choice but in reverse order.
  
\item Parties occasionally chose to vary the method of sharing
  treasure but any such variation should be discussed prior to leaving
  the guild and included in the adventuring contract signed by all
  party members.
\end{itemize}

\subsubsection{Debriefing}

When each guild adventure returns Guild Security requires the
adventurers to describe and justify their actions especially with
reference to the Guild’s reputation, the law and the membership
contract.

During this debriefing the Guild will also check if the member is
controlled by outside agencies, has been cursed or requires
healing. They will discuss with the member the results and the best
way of dealing with any such effects.

\subsubsection{Guild Taxes}

The Guild collects from members 10\% of all gross income generated on
adventure. Part of this money is passed onto the Duke in payment of
Carzalan taxes for members; the remainder is used to pay for the
services provided by the Guild.

\subsection{Getting a Rulebook}

While our campaign is primarily based on DragonQuest™ (2nd ed.), the
GMs have made many changes to the rules over the years. We reissue a
rulebook in loose-leaf format every so often, and this is when rules
changes come into play. This is available from the DQ Wiki (see below)
as a pdf.

\subsection{Web Pages}

https://dq-nz.org/dqwiki/index.php/Main\_Page

\subsection{Mailing Lists}

dqnz-announce@googlegroups.com

\subsection{Other Documents}

There are a few resource books which are available on the wiki for
both players and GMs. Some of these are:

\begin{itemize}
\item DQ Rulebook 2021
\item A guide to the Adventures Guild of Seagate – June 2021 (this
  document)
\item The Official DQ GMs Guide – Dec 2005 (a document aimed at
  assisting GMs)
\item The Duchy of Carzala. (36 page Source booklet detailing Carzala)
\end{itemize}

\subsection{GM Contact List}

The following GMs are happy to be contacted for questions, character
generation etc.  Contact by email is generally preferable.  Remember,
if telephoning, to choose a reasonable hour.

\begin{tblr}{colspec={ll}}
Stephen & sm1971@gmail.com \\
Jonathan & jonobean@gmail.com \\
Keith & phaetonnz@gmail.com \\
\end{tblr}

\subsection{New GMs welcome}

We always need new GMs joining the DQ campaign. When dealing with new
GMs please try to support them as best you can.  Players have the
largest impact on new GMs and with your effort and support we can
encourage and retain new GMs.

If you have thought about running a game, give it a go. Talk to one of
the GMs from the GM contact list.

We can offer a wide range of support for starting GMs ranging from:

\begin{itemize}
\item Help with Alusia World information.
\item Storyboarding \& adding timelines to your adventure, to help fit
  into the DQ season.
\item Offering advice \& experience with balancing encounters.
\item Help with writing up items.
\end{itemize}

If you are interested in either GMing or if you need help from another
GMs, please let Stephen.

\subsection{Fair Play \& Courtesy}

Our campaign flavour is that of heroic fantasy and Guild members are
the good guys who stick together through thick and thin. However being
a good guy all the time can get a little dull, and players sometimes
like to portray a rogue. We have a long running campaign and therefore
unlike some other campaigns we have some artificial rules laid out.

\begin{itemize}
\item Don't cheat other characters.
\item Temper your character's behaviour to ensure fun for everyone. We
  have a co-operative role-playing setting, and not a competitive one.
\item Be nice to the Game Master - they are the people who hand out
  the experience points. They are working hard to make a fun game for
  all.  Please try to make sure that they are having a good time too.
\item Be nice to the other players - We only sign up for games every
  13 weeks so it is along time between meetings.
\item Have fun - if you aren’t having fun remember to eat another
  chocolate biscuit, sugar high always brings out the best in tense
  role-playing situations.
\end{itemize}

Guild Security is a mechanism of wide and undefined powers to ensure
fairness amongst the players, at any time the Game Master might feel
the need to use guild security to step in to either mediate or
“punish” a rogue character.

\begin{itemize}
\item If you have a disagreement with a GM over a rule, or
  interpretation, say so.  GMs are not always right, but please say it
  only once. If the GM stands by the ruling, stop arguing and get on
  with the game. If it really is important bring it up again at the
  end of the session. If you still can’t settle it, talk to some of
  the other GMs and a solution will be found.

\item If you can’t make it to a game session let the GM know so they
  don’t waste time waiting for you. It is rude and inconsiderate to
  keep people waiting. If you can’t turn up to most of the sessions of
  a game, or can’t be bothered, don’t go on the adventure.  Most GMs
  balance their game for the number of players booked and it is
  extremely frustrating to have it ruined by players who don’t turn
  up.

\item Don’t intimidate other players or play their characters for
  them.  If you can’t find a way to give advice in character, then
  don’t.  If someone really needs help, offer it but don’t force it on
  them.  It’s their right to make mistakes too.
\end{itemize}

\subsection{Grievance Tribunal}

The Grievance Tribunal is to settle disputes. If you feel
uncomfortable with something your GM is doing or you feel you have
been unfairly treated by your GM or another player you should
initially discuss this with your GM outside of the game. If you can’t
resolve the problem with them, then you may take your grievances to
the tribunal. The tribunal represents all GMs and will adjudicate on
the matter. This should be seen as an option of last resort.

The current tribunal members are listed on DQ Wiki (see Web Pages)

\subsection{Character Generation}

To play in our DQ campaign, you will have to generate a new character
with one of the GMs. It is not allowed for players to “move”
characters into DQ from other campaigns or systems.

Before you contact the GM you should think about what type of
character you would like to play. After you have the basic idea,
approach one of the GMs and gain their assistance to generate the
character.

Although you may generate a character completely randomly, most
aspects of a character may be chosen. The most fundamental choices are
those of race and whether or not to be a magic user. Ask your friends
or GM to précis the races and magic colleges for you if you are
unsure, but be aware that, no matter how unbiased they try to be, any
GM will have more experience of some areas than others and hence
cannot help but be a little biassed.

An important aspect of the process of character generation is the
development of a defined persona for your character. Players are
encouraged to develop a history and a set of basic motivations for
their character.  Characters are assumed to have spent six months
training with the Guild before being admitted as members. This covers
learning a college or the basic fighting and survival skills.  Think
about how was your training paid for?  Did your family help; did you
clean the Guild stables for 3 years, or some other arrangement?

You can have more than one character in DQ, but you are asked to play
in only one game per session (otherwise some other player will miss
out on playing). We recommend that you start a second character no
sooner than six months after starting play. In this time you may gain
a better understanding of the system, and what type of character you
would like to play next.

\subsection{The “In-Game” Guild}

The Seagate Adventurers Guild is in the Duchy of Carzala and is
located five minutes south of the city of Seagate on the other side of
the Sweetwater River. It is a walled and guarded enclave.

The Guild enclave contains:
\begin{itemize}
\item Guild Meeting Hall and attached small common rooms (normally
  used for party and employer meetings etc).

\item The Library is in the new Council House. It is split into
  sections including the Guild records, Adventure scribe notes,
  Historical, Geographical, Legal, Social and Cartography.  Guild
  Security and the Vaults. Guild Security have a public office in the
  Council House but also have guard rooms on the gates and in other
  buildings. The Vaults are accessed though the Council House cellar.

\item Residential area, including sleeping quarters dining hall,
  kitchen, storerooms and bathhouses.

\item The magic departments spread across several buildings.

\item Casting Chambers for magical training. These are soundproof and
  regularly checked when in use.

\item Military Department including offices, armoury, smithy, and
  arena.

\item Hospital.

\item Guild Pub, rather than a main common room.

\item Stables and stable hand’s quarters.

\end{itemize}

The Guild provides a common meeting ground for adventurers, to enable
training and research in relative safety and a pool of information,
resulting in a close knit and more effective band of Adventurers.

The main services provided to members include: diverse magical
services; healing and resurrection; valuation and clearing house for
treasure; banking and strong rooms; a library; training facilities in
magic, healing and military skills; co-ordination for other training;
and legal services.

\end{multicols}
